\subsection{不变双线性形式}

设 g 是 K 上的李代数。g 在 g 上的伴随表示以及 g 在 K 上的零表示，在 K-模

\(N = \mathcal{L}(g, g; K)\)（即 g 上的双线性形式组成的 K-模）上定义了 g-模结构。简言之，若 g 上的双线性形式 \(\beta\) 在表示 \(x \mapsto x_{N}\) 下不变，则称其为不变的。由公式 (10)，这成立的充要条件是：

\[
\beta ([ x, y ], z) = \beta (x, [ y, z ])\tag{13}
\]

对所有 \(x, y, z\) 属于 \(\mathfrak{g}\) 成立。

现在设 \(\mathfrak{d}\) 是 \(\mathfrak{g}\) 的导子李代数。\(\mathfrak{d}\) 的恒等表示以及 \(\mathfrak{d}\) 在 \(\mathrm{K}\) 上的零表示，定义了 \(\mathrm{D} \mapsto \mathrm{D}_{\mathrm{N}}\) 在 \(\mathfrak{d}\) 上的表示 \(\mathrm{N}\)。简言之，若 \(\mathfrak{g}\) 上的双线性形式在表示 \(\mathrm{D} \mapsto \mathrm{D}_{\mathrm{N}}\) 下不变，则称其为完全不变的。完全不变的双线性形式是不变的。对于 \(\beta\) 上的双线性形式 \(\mathfrak{g}\)，完全不变的充要条件是：

\[
\beta (\mathrm{D} x, y) + \beta (x, \mathrm{D} y) = 0\tag{14}
\]

对所有 \(x, y\) 属于 \(\mathfrak{g}\) 且 \(\mathrm{D} \in \mathfrak{d}\) 成立。

命题 7. 设 \(\mathfrak{g}\) 是李代数，\(\beta\) 是 \(\mathfrak{g}\) 上的不变对称双线性形式，\(\mathfrak{a}\) 是 \(\mathfrak{g}\) 的理想。

\begin{enumerate}
\def\labelenumi{(\alph{enumi})}
\item
  \(\mathfrak{a}'\) 关于 \(\mathfrak{a}\) 的正交子空间 \(\beta\) 是 \(\mathfrak{g}\) 的理想。
\item
  如果 \(\mathfrak{a}\) 是特征理想且 \(\beta\) 完全不变，则 \(\mathfrak{a}'\) 是特征理想。
\item
  如果 \(\beta\) 非退化，则 \(\mathfrak{a} \cap \mathfrak{a}'\) 是交换的。
\end{enumerate}

设 D 是 g 的一个导子。假设 a 在 D 下稳定，并且对 x,y 属于 g 有 \(\beta(\mathrm{D}x,y)+\beta(x,\mathrm{D}y)=0\)。则 \(z\in a'\) 蕴含 \(Dz\in a'\)，因为对所有 \(t\in a\)，都有 \(Dt\in a\)，从而 \(\beta(\mathrm{D}z,t)=-\beta(z,\mathrm{D}t)=0\)。因此 \(a'\) 在 D 下稳定。这就证明了 (a) 和 (b)。

现在设 \(\mathfrak{b}\) 是 \(\mathfrak{g}\) 的理想，并假设 \(\beta\) 在 \(\mathfrak{b}\) 上的限制为零。对 \(x, y\) 属于 \(\mathfrak{b}\) 且 \(z \in \mathfrak{g}\)，有 \(\beta ([ x, y ], z) = \beta(x, [y, z]) = 0\)，因为 \([y, z] \in \mathfrak{b}\)。因此 \([\mathfrak{b}, \mathfrak{b}]\) 与 \(\mathfrak{g}\) 正交。若 \(\beta\) 非退化，则 \(\mathfrak{b}\) 必为交换的。将此结果应用于 \(\mathfrak{a} \cap \mathfrak{a}'\)，即得 (c)。

定义 4. 设 g 是 李 K-代数，M 是 g-模。假设把 M 看作 K-模时有有限基。与 g-模 M（或相应表示）关联的双线性形式，是 g 上的对称双线性形式 \((x,y)\mapsto\mathrm{Tr}(x_{\mathrm{M}} y_{\mathrm{M}})\)。如果所考虑的表示是伴随表示，则相应的双线性形式称为 g 的 Killing 形式。

命题 8。设 \(\mathfrak{g}\) 是李代数，\(\mathbf{M}\) 是 \(\mathfrak{g}\)-模。假设把 \(\mathbf{M}\) 看作 K-模时有有限基。与 \(\mathbf{M}\) 关联的双线性形式是不变的。对 \(x, y, z\) 属于 \(\mathfrak{g}\)，有：

\[
\begin{array}{r l} & {\mathrm{Tr} ([ x, y ] _ {\mathrm{M}} z _ {\mathrm{M}}) = \mathrm{Tr} (x _ {\mathrm{M}} y _ {\mathrm{M}} z _ {\mathrm{M}}) - \mathrm{Tr} (y _ {\mathrm{M}} x _ {\mathrm{M}} z _ {\mathrm{M}}) = \mathrm{Tr} (x _ {\mathrm{M}} y _ {\mathrm{M}} z _ {\mathrm{M}}) - \mathrm{Tr} (x _ {\mathrm{M}} z _ {\mathrm{M}} y _ {\mathrm{M}})} \\ & {\quad = \mathrm{Tr} (x _ {\mathrm{M}} [ y, z ] _ {\mathrm{M}}).} \end{array}
\]

命题 9. 假设 K 是域，李代数 g 在 K 上有限维。设 a 是 g 的理想，β 是 g 的 Killing 形式，β' 是 a 的 Killing 形式。则 β' 是 β 在 a 上的限制。

设 u 是向量空间 g 的一个使 a 保持不变的自同态。设 v 是 u 在 a 上的限制，w 是 u 在取商后诱导的向量空间 g/a 的自同态。取 g 的一个基 \(Tr u = Tr v + Tr w\)，其中前 p 个元素构成 a 的基，便可看出 \((x_{1}, \ldots, x_{n})\)。现在令 \(x \in a, y \in a\)，将上述公式应用于 \(u = (\mathrm{ad}_{\mathrm{g}} x)(\mathrm{ad}_{\mathrm{g}} y)\) 的情形。此时 \(v = (\mathrm{ad}_{\mathrm{a}} x)(\mathrm{ad}_{\mathrm{a}} y)\)，且 w = 0。因此 \(\beta(x, y) = \beta'(x, y)\)。

命题 10. 假设 K 是域，李代数 g 在 K 上有限维。g 的 Killing 形式 β 是完全不变的。

设 D 是 g 的导子。存在一个李代数 \(g'\)，其中 g 是余维 1 的理想，并且存在 \(x_{0}\) 中的元素 \(g'\)，使得对所有 \(Dx = [x_{0}, x]\) 都有 \(x \in g\)（§ 1，第 8 号，例 1）。设 \(\beta'\) 是 \(g'\) 的 Killing 形式。对 g 中的 x、y，有 \(\beta'([x, x_{0}], y) = \beta'(x, [x_{0}, y])\)，即 \(\beta'(Dx, y) + \beta'(x, Dy) = 0\)。现在 \(\beta'\) 在 g 上的限制是 \(\beta\)（命题 9）。命题得证。

